% =======================================================================
% pw_temporal_genesis_end_to_end_revised.tex
% From No Clock to Relational Time (PW) -> Engine -> AR(inf)->AR(2) -> Decoherence
% Clock births, multiplicities, GR bridge; operational PBC/APBC; directed graphs; reproducibility
% =======================================================================
\documentclass[11pt]{article}

\usepackage[margin=1in]{geometry}
\usepackage{amsmath,amssymb,mathtools,bm}
\usepackage{physics}
\usepackage{graphicx}
\usepackage{booktabs}
\usepackage{lmodern}
\usepackage{microtype}
\usepackage{hyperref}
\usepackage[most]{tcolorbox}
\usepackage{xcolor}
\usepackage{tikz}
\usetikzlibrary{arrows.meta,calc,positioning}
\hypersetup{colorlinks=true,linkcolor=blue,citecolor=blue,urlcolor=blue}

% --- Box styles ---
\tcbset{sharp corners, colback=white, boxsep=1.0ex}
\newtcolorbox{mydef}[1]{breakable,colframe=black!55,colback=black!2,title=\textbf{#1}}
\newtcolorbox{mythm}[1]{breakable,colframe=blue!60!black,colback=blue!2,title=\textbf{#1}}
\newtcolorbox{myprop}[1]{breakable,colframe=green!60!black,colback=green!2,title=\textbf{#1}}
\newtcolorbox{mybox}[1]{breakable,colframe=black!40,colback=black!3,title=\textbf{#1}}
\newtcolorbox{proofbox}[1]{breakable,colframe=black!60,colback=black!1,title=\textbf{#1}}

% --- Shortcuts ---
\newcommand{\Id}{\mathsf{I}}
\newcommand{\Adj}{\mathsf{A}}
\newcommand{\Lap}{\mathsf{L}_\alpha}
\newcommand{\Thet}{\Theta}
\newcommand{\M}{M}
\newcommand{\V}{V}
\newcommand{\Vl}{V_\ell}
\newcommand{\Q}{Q}
\newcommand{\Qell}{Q_\ell}
\newcommand{\plz}{p_\ell(z)}
\newcommand{\Hol}{\mathrm{Hol}}
\newcommand{\LC}{\mathrm{LC}}
\newcommand{\E}{\mathsf{E}}
\newcommand{\sgn}{\mathrm{sgn}}

\title{\vspace{-0.6cm}\Large From No Clock to Relational Time:\\
\large Page--Wooters Genesis, AR($\infty$)$\to$AR(2), Clock Births, Multiplicities, and a Finite-Loop GR Bridge}
\author{}
\date{\today}

\begin{document}
\maketitle
\vspace{-1.2em}\hrule\vspace{0.6em}

\begin{abstract}
We reformulate the program ``from no intrinsic clock to emergent clocks, multiplicities, and geometry'' in a relational (Page--Wooters) setting. Layer~0 is a \emph{clockless label space} on which regulators and probes act; temporal order appears only after a gauge choice. The conditioned dynamics of residue streams obey an exact decimated law derived from a $2\times2$ transfer engine. We show how an AR($\infty$) parent with one dominant conjugate pair admits a controlled AR(2) surrogate, and how decoherence arises as phase diffusion quantified by $D=\tfrac12 S_\delta(0)$ under demodulation. Clock \emph{births} are produced by net-zero probes with closed-form (finite-ring) Green functions; a binary PBC/APBC probe opens/closes a clock-sector gap and flips an \emph{operational} trace estimator $\widehat V_L$ while the intrinsic $V_L=\mathrm{tr}(M^L)$ remains invariant. Finally, we provide a finite-loop holonomy--Levi--Civita bridge with explicit units and a clear error budget, extend to directed graphs, and include code-ready estimators.
\end{abstract}

\tableofcontents
\vspace{0.6em}\hrule\vspace{0.8em}

% =======================================================================
\section{Layer 0: No intrinsic time --- labels, regulators, and probes}
% =======================================================================
\noindent At the foundational layer we assume only a set of labels $\mathcal{E}$ (no metric, order, or topology). \emph{Regulators} are symmetry operations $R\in\mathcal{R}$ acting on functions $f:\mathcal{E}\!\to\!\mathbb{C}^2$ (e.g., group representations, quotientings), while \emph{probes} are linear functionals extracting finite summaries. No sequential structure is postulated. A \emph{gauge choice} is a section that identifies a discrete cover of a regulator orbit with $\mathbb{Z}$, yielding an index map $k\!\mapsto\!e_k$ (bookkeeping). Only after this choice do recursions or boundary conditions appear; different gauges are physically equivalent if they produce indistinguishable probe statistics.

\paragraph{Resolution of the ``No Clock'' paradox.}
Objects such as Möbius twists, difference operators, or Dirac boundary conditions \emph{do not} assert pre-existing time; they are \emph{regulators} and \emph{diagnostics} on label classes. Temporal order emerges when a gauge is \emph{fixed} and validated by acceptance tests (e.g., stability of inferred invariants across gauges).

% =======================================================================
\section{Relational conditioning (Page--Wooters) in discrete form}
% =======================================================================
Let the composite ``clock $\otimes$ system'' be constrained by a discrete PW equation $\widehat{\mathcal{C}}\ket{\Psi}=0$ with $\widehat{\mathcal{C}}$ block-linear in a shift on the clock labels and a transfer on the system. Conditioning on a clock outcome (a probe event) and choosing a gauge section $\mathbb{Z}\!\to\!\mathcal{E}$ yields linear residue dynamics on the system observables $\{x_k\}$ with a \emph{transfer engine} described next.

% =======================================================================
\section{Engine: GL$(2,\mathbb{C})$ transfer, exact decimation, and speed bound}
% =======================================================================

\begin{mydef}{Transfer engine and invariants}
A one-step transfer $\M\in\mathrm{GL}(2,\mathbb{C})$ (gain/loss allowed). Define
\[
\V_1=\mathrm{tr}\,\M,\qquad \Q=\det\M,\qquad \Vl=\mathrm{tr}(\M^\ell),\qquad \Qell=(\det \M)^\ell.
\]
For a scalar observable $\{x_k\}$, decompose into residue tracks $y^{(r)}_t=x_{r+\ell t}$.
\end{mydef}

\begin{mythm}{Exact decimated two-step law (Cayley--Hamilton)}
For each residue $r$ and stride $\ell$,
\begin{equation}
y^{(r)}_{t+2}-\Vl\,y^{(r)}_{t+1}+\Qell\,y^{(r)}_{t}=0,\qquad
r_\pm(\ell)=\frac{\Vl\pm\sqrt{\Vl^2-4\Qell}}{2}.
\label{eq:engine-law}
\end{equation}
\end{mythm}

\begin{proofbox}{Sketch}
Cayley--Hamilton on $M^\ell$: $M^{2\ell}-V_\ell M^\ell + Q^\ell \Id=0$.
Apply to state $\Psi_k$ and project to any linear observable.
\end{proofbox}

\paragraph{Unitary vs transfer engine.}
If $M\in\mathrm{SU}(2)$ then $\Q=1$ and $|r_\pm|=1$. For general $M\in\mathrm{GL}(2,\mathbb{C})$, $\Q\neq1$ captures amplification/attenuation; one may optionally normalize $\widetilde M:=M/\sqrt{\det M}\in\mathrm{SL}(2,\mathbb{C})$ for trace identities.

\paragraph{Envelope and speed limit (with units).}
Introduce a locality envelope
\begin{equation}
\E := \exp\!\big(-\sigma^2\kappa\,\tau\,\Lap^\dagger\Lap\big),\qquad s:=\sigma^2\kappa\,\tau\ \ (\text{dimensionless}),
\end{equation}
with $\Lap$ a (possibly PSD) graph Laplacian, $\Lambda_{\max}$ its spectral radius, and $\tau$ the step.
Then
\begin{equation}
v_{\max}\ \lesssim\ \frac{a\,\sup_{\mathbf q}\|\nabla_{\mathbf q}\varepsilon(\mathbf q)\|}{1+s\,\Lambda_{\max}^2}
\quad\text{(Fourier band phase $\varepsilon$)},\qquad
v_{\max}\ \lesssim\ \frac{2 a\,\|\widehat K\|\,\Delta}{1+s\,\Lambda_{\max}^2}\quad\text{(graph)}.
\label{eq:speed}
\end{equation}

% =======================================================================
\section{AR($\infty$) parent $\to$ AR(2) surrogate and consistency regime}
% =======================================================================
Assume a coarse observable has an AR($\infty$) parent with one dominant conjugate pair $r_\pm=\rho e^{\pm i\theta}$ ($\rho\le1$) and a stable tail.

\begin{myprop}{Padé [1/1] surrogate captures the pair exactly}
The AR(2) surrogate
\begin{equation}
x^{(2)}_{t+2}-V\,x^{(2)}_{t+1}+Q\,x^{(2)}_{t}=u_t,\qquad V=2\rho\cos\theta,\quad Q=\rho^2,
\label{eq:ar2-surrogate}
\end{equation}
has characteristic roots $r_\pm=\rho e^{\pm i\theta}$; the tail folds into an MA($\infty$) innovation $u_t$.
\end{myprop}

\paragraph{When is AR(2) adequate?}
Empirically require a dominance ratio
\[
\frac{|r_+|}{\max_{j\ge 3}|p_j|}\ \gtrsim\ 1.2\quad\text{and white AR(2) residuals (Ljung--Box at 5\%).}
\]
Specialize to $Q=1$ (unitary engine) by setting $\rho=1$ and $V=2\cos\theta$.

\begin{figure}[h]
\centering
\includegraphics[width=0.6\linewidth]{ar2_pade_sweep.png}
\caption{AR($\infty$)$\to$AR(2) surrogate sweep: fit quality across tail strengths and dominance ratios.}
\end{figure}

% =======================================================================
\section{Superposition and decoherence as phase diffusion}
% =======================================================================
Project onto the rotating eigenmode, demodulate by the fitted $\hat r_+$, and write the phasor recursion
\begin{equation}
a_{t+1}=\rho\,e^{i(\theta+\delta_t)}\,a_t,\qquad \mathbb{E}\,\delta_t=0.
\end{equation}
Then off-diagonal coherence decays as
\begin{equation}
|\mathcal{C}(t)|:=\frac{|\mathbb E[a_t\overline{a_0}]|}{\mathbb E|a_0|^2}\ \approx\ \rho^{\,t}\,e^{-\tfrac12 S_\delta(0)\,t},
\qquad D:=\tfrac12 S_\delta(0).
\label{eq:dephasing}
\end{equation}
\emph{From additive noise to phase noise.} If $a_{t+1}=r_+ a_t+\xi_t$ with small complex innovation, then
$\delta_t\approx \Im\!\big(\xi_t/(r_+ a_t)\big)$; estimate $S_\delta(0)$ from the demodulated, amplitude-normalized residual.

\begin{figure}[h]
\centering
\includegraphics[width=0.6\linewidth]{coherence_fit.png}
\caption{Decoherence fit $|\mathcal C(t)|$ vs the model $\rho^t e^{-Dt}$ with $D=\tfrac12 S_\delta(0)$.}
\end{figure}

\begin{mybox}{One-pass estimator for $(\rho,D)$ (pseudocode)}
\small
\begin{verbatim}
Fit AR(2): y_{t+2} = V y_{t+1} - Q y_t + u_t  ->  roots r_+/-, pick r_+ (arg in (0, pi))
Demodulate: z_t = y_t * (conj(r_+))^t
Phase increments: delta_t = arg(z_{t+1}*conj(z_t)) - arg(r_+)  (wrapped to (-pi, pi])
Estimate S_delta(0) via Bartlett at zero; return rho=|r_+|, D = 0.5*S_delta(0)
\end{verbatim}
\end{mybox}
\noindent\emph{Robustness.} Gate out very small $|a_t|$ or weight by $|a_t|^2$ when estimating $S_\delta(0)$ to stabilize phase-noise estimates at low amplitudes.

\paragraph{PBC/APBC and $D$.}
A PBC (Ramond) clock admits a zero mode; APBC (NS) opens a gap $\sim\pi/T$. This \emph{modifies} the low-frequency phase-noise power $S_\delta(0)$ through clock--tail coupling (sign model-dependent); test empirically by the above estimator. See also Sec.~\ref{sec:binary}.

% =======================================================================
\section{Gravity as AR($\infty$) Back–Reaction of Clocks (Newtonian Limit)}
% =======================================================================
\label{sec:gravity-arinfty}

\paragraph{From AR($\infty$) to a static slowness field.}
A coarse observable admits a two–pole + tail representation
\[
x_{t+1}-Vx_t+Qx_{t-1}=\sum_{j\ge1} b_j x_{t-j}+u_t,\qquad
V=2\rho\cos\theta,\ \ Q=\rho^2.
\]
Demodulate by the dominant pole \(r_+=\rho e^{i\theta}\): \(z_t=x_t r_+^{-t}\).
The tail induces phase diffusion \(z_{t+1}=e^{i(\delta_t-\phi(x))}z_t\) with zero–mean \(\delta_t\) and a slow, static
\(\phi(x)\) (phase–slowness field). Define the source density
\[
\rho_{\text{tail}}(x)=\beta\,S_\delta(x,0),
\]
with \(S_\delta(x,0)\) the zero–frequency power of phase increments (estimated after AR(2) demodulation).

\paragraph{Envelope as a gradient penalty.}
The envelope \(E=\exp(-s\,\Lap^\dagger\Lap)\) with \(s=\sigma^2\kappa\,\tau\) penalizes spatial phase gradients, yielding the static functional
\[
\mathcal J[\phi]=\frac{\kappa_s}{2}\int |\nabla\phi|^2\,dV-\int \phi\,\rho_{\text{tail}}\,dV,\qquad
\kappa_s:=s\,c_{\mathrm{env}}.
\]
The Euler–Lagrange equation gives \(-\kappa_s\nabla^2\phi=\rho_{\text{tail}}\).
With \(\Phi:=\alpha\phi\) and \(4\pi G:=\alpha^{-1}\kappa_s^{-1}\beta\), we obtain the Poisson equation
\[
\boxed{~\nabla^2\Phi=4\pi G\,\rho~,}\qquad \rho:=\rho_{\text{tail}}.
\]

\paragraph{Interpretation.}
\(\Phi\) is the emergent Newtonian potential (phase–time dilation);
\(\rho\) is the density of low–frequency tail power (or clock–birth intensity).
In the pure AR(2) limit (\(\delta_t\equiv0\)) we have \(\rho=0\Rightarrow \nabla^2\Phi=0\).
Hence gravity is an AR($\infty$) effect, not an AR(2) effect.

\paragraph{Consistency with the binary probe.}
The PBC/APBC dial modifies access to \(\omega\approx 0\) in the clock sector (zero mode vs gap),
changing \(S_\delta(0)\) and thus \(\rho\) operationally, while the intrinsic system trace \(V_\ell=\mathrm{tr}(M^\ell)\) is unchanged
(the observed estimator \(\widehat V_\ell\) flips sign after conditioning under APBC).
% (Background: decimated law, Green functions, and the PBC vs APBC spectrum for the time circle.)

% =======================================================================
\section{Clock births by net-zero probes and finite-ring Green functions}
% =======================================================================
Inject a net-zero source $s$ on a macrocycle of length $L$ (e.g.\ $s_m=+\delta_{m,0}-\delta_{m,L/2}$ for even $L$). The AR(2) response $x$ solves
\begin{equation}
x_{m+1}-V x_m + Q x_{m-1}=s_m,\qquad \sum_{m=0}^{L-1} s_m=0.
\end{equation}
\paragraph{Finite ring Green function (enforcing net-zero).}
On a ring,
\begin{equation}
G^{(L)}_n=\frac{1}{L}\sum_{k=0}^{L-1}\frac{e^{2\pi i k n/L}}{V-2\sqrt{Q}\cos(2\pi k/L)},
\label{eq:ring-G}
\end{equation}
reducing to $G_n=\sin(n\phi)/\sin\phi$ for $Q=1$, $V=2\cos\phi$ as $L\to\infty$.
\emph{Net-zero source:} the DC Fourier bin ($k=0$) is not excited, so there is no DC response.
Handle the resonant case $V=2\sqrt{Q}$ by the double-root limit/pseudoinverse.

The resulting response $x=G^{(L)}\!*\,s$ exhibits clean ``births'' of oscillatory residues shaped by the envelope.

\begin{figure}[h]
\centering
\includegraphics[width=0.8\linewidth]{green_birth.png}
\caption{Clock births from a net-zero source on one macrocycle and repeated on the universal cover.}
\end{figure}

% =======================================================================
\section{Binary probe: PBC zero mode vs APBC gap (operational flip)}
\label{sec:binary}
% =======================================================================
A discrete Dirac on a time circle of period $T$ has eigenfrequencies
$\omega_n=\tfrac{2\pi}{T}(n+\tfrac{s}{2})$ for $s=0$ (PBC) or $s=1$ (APBC):
PBC admits a zero mode; APBC has $|\omega|\ge \pi/T$ (gap). In our operational pipeline, twisting the \emph{clock} by APBC injects a $\pi$-phase over one macrocycle in the conditioned residue at stride $L$, which flips the \emph{estimator} $\widehat V_L$ derived from the data, while the intrinsic $V_L=\mathrm{tr}(M^L)$ is unchanged (a system invariant). This matches the spectral picture and the checksum demo reported in the supporting note. \emph{Caption fix for plots:} “$\widehat V_L$ flips sign under APBC; $V_L$ is invariant.”  % (Binary holonomy test and spectral proof)

\begin{figure}[h]
\centering
\includegraphics[width=0.75\linewidth]{dirac_pbc_apbc_gap.png}
\caption{Discrete Dirac on a time circle: PBC shows a zero mode; APBC is gapped by $\sim\pi/T$.}
\end{figure}

\begin{figure}[h]
\centering
\includegraphics[width=0.8\linewidth]{binary_holonomy_signflip.png}
\caption{Operational test: $\widehat V_L$ flips sign under APBC while $V_L$ remains invariant.}
\end{figure}
% (External note with PBC zero vs APBC gap and checksum example)
% [Reference note acknowledged in the covering letter]
% (In chat: see companion citation)

\begin{proofbox}{Least-squares sign flip at stride $L$}
The AR(2) LS estimator on the stride-$L$ residue is
$\displaystyle \widehat V_L=\frac{\sum_t (y_{t+2}\overline{y_{t+1}}+y_t\overline{y_{t+1}})}{\sum_t |y_{t+1}|^2}$.
APBC sends $y_t\mapsto(-1)^t y_t$. The numerator gains a factor $(-1)$ while the denominator
is unchanged; hence $\widehat V_L\mapsto-\,\widehat V_L$ (and $\widehat Q_L$ is unchanged).
\end{proofbox}

% =======================================================================
\section{Multiplicity and species: operational labels and how to measure them}
% =======================================================================
We label species by operational invariants
\[
\Lambda=(\ell,\ V_\ell,\ Q_\ell,\ r_\pm=\rho e^{\pm i\theta},\ \chi\in\{\pm1\},\ (c_0,c_1),\ \mathrm{hol}),
\]
where $\chi$ is the PBC/APBC parity measured at $L=\ell$, $(c_0,c_1)$ are stream checksums (TKNN-like integers), and ``hol'' is the holonomy class (circle vs Möbius/Pin$^\pm$).
\paragraph{How to compute $(c_0,c_1)$ quickly.}
Over a $k$-window, accumulate two linear invariants (real/imag parts) of $\sum_k d_k r_+^k$; tri-site local moves preserve them. A soft penalty with weight $\nu$ stabilizes reconstruction; normal equations $(I+\nu A^\top A)d=m+\nu A^\top b$ yield digits $d$ (we reserve $\mu$ for the spectral gap of $L_\alpha^\dagger L_\alpha$).  % operational summary; detailed formulae in Appendix

\paragraph{Independence scale.}
If $L_\alpha^\dagger L_\alpha$ has spectral gap $\mu>0$, then $\mathrm{Cov}\!\big(y^{(r)},y^{(r+d)}\big)\sim e^{-\mu d}$.
In practice take $d\gtrsim 3\,\mu^{-1}$ to treat residue tracks as independent samples.

% =======================================================================
\section{Finite-loop GR bridge with units, assumptions, and interpretation}
% =======================================================================
Let phases $\varphi^a$ push forward to a coframe $e^a=\Thet^a{}_i\,dx^i$; define $F^a{}_{ij}=\partial_i\Thet^a{}_j-\partial_j\Thet^a{}_i$, the Levi--Civita connection $\omega$, and curvature $\Omega$.
Compare envelope holonomy $P_{\mathrm{env}}(C)$ with Levi--Civita $H_{\LC}(C)$ around loop $C$ with spanning surface $\Sigma(C)$:
\begin{mybox}{Finite-loop bridge (error budget with units)}
\begin{equation}
\big\|P_{\mathrm{env}}(C)-H_{\LC}(C)\big\|
\;\le\;
C\Big(
\underbrace{\sqrt{\int_{\Sigma(C)}\!\!\|F\|^2\,dA}}_{\text{antisym.\ defect}}
+\underbrace{h}_{\text{mesh}}
+\underbrace{\ell^3\|\nabla\Omega\|_\infty}_{\text{curv.\ var.}}
+\underbrace{\mu^{-1}\,s\,\ell}_{\text{envelope leakage}}
\Big),
\label{eq:bridge}
\end{equation}
\end{mybox}
with $s=\sigma^2\kappa\,\tau$, $\mu$ the spectral gap (units $L^{-2}$) of $\Lap^\dagger\Lap$, loop diameter $\ell$, and mesh $h$. Norms are operator norms on $\mathfrak{so}(d)$; constants absorb unit conversions.
\paragraph{Idea of proof.} (i) Strang splitting + BCH on edges (commutator budget $\propto s$); (ii) non-Abelian Stokes to replace edge-ordered exponentials by a surface integral (curvature budget); (iii) Cartan + Stokes + Poincaré/trace on rectangles, then polygonal refinement (defect + mesh budgets). 

\begin{figure}[h]
\centering
\includegraphics[width=0.65\linewidth]{bridge_scaling.png}
\caption{Finite-loop bridge calibration: scaling of $\|P_{\mathrm{env}}-H_{\LC}\|$ against the four error-budget terms.}
\end{figure}

% =======================================================================
\section{Directed/weighted graphs (safe PSD Laplacian) and scaling knobs}
% =======================================================================
For directed or weighted $\Adj$, use the PSD surrogate
\[
\mathsf{L}_{\mathrm{psd}}=(\Id-\alpha\Adj)^\dagger(\Id-\alpha\Adj),\qquad 0<\alpha<\|\Adj\|_2^{-1},
\]
and replace spectral radii by singular values throughout. The speed and bridge bounds retain their form with constants adjusted.
\paragraph{Knobs $\to$ effects (quick map).}
\begin{center}
\begin{tabular}{@{}lll@{}}
\toprule
Knob & Primary effect & Practical setting \\
\midrule
$s=\sigma^2\kappa\,\tau$ & locality $\uparrow$, $v_{\max}\downarrow$ & target $s\,\Lambda_{\max}^2\simeq 1$ for 50\% speed cut \\
$\gamma$ (curl penalty) & $\|F\|$ suppression & raise until $\sqrt{\int\|F\|^2}\le 0.01\,\mathrm{Area}$ \\
$\beta$ (pushforward tie) & enforce $\Thet\approx\nabla\varphi$ & set $\beta\approx \gamma a^2$ \\
APBC twist & gap/zero-mode toggle & A/B test $\widehat V_L$, measure $\Delta D$ \\
\bottomrule
\end{tabular}
\end{center}

% =======================================================================
\section{Minimal experiments (falsifiable), figures, and captions}
% =======================================================================
\begin{itemize}
\item \textbf{PBC vs APBC spectrum (clock sector):} show zero mode vs gap on the time circle (discrete Dirac); caption: “PBC has a zero mode; APBC opens a gap $\sim\pi/T$.”
\item \textbf{Operational flip:} plot $\widehat V_L$ from the conditioned residue under PBC vs APBC; caption: “$\widehat V_L$ flips sign under APBC; $V_L$ is invariant.”
\item \textbf{AR($\infty$)$\to$AR(2) adequacy:} sweep tails; report residual whiteness (LB test) and MSE vs dominance ratio.
\item \textbf{Decoherence:} estimate $(\rho,D)$ from demodulated phase; overlay $|\mathcal{C}(t)|$ vs $\rho^t e^{-Dt}$.
\item \textbf{Bridge calibration:} vary $(s,h,\ell)$ on a grid with known $\Omega$; verify slopes in \eqref{eq:bridge}.
\end{itemize}

% =======================================================================
\section{Symbols and units (first-use table)}
% =======================================================================
\begin{center}
\begin{tabular}{@{}lll@{}}
\toprule
Symbol & Units & Meaning \\
\midrule
$a$ & length & lattice spacing \\
$\tau$ & time & step duration \\
$s=\sigma^2\kappa\,\tau$ & 1 & envelope strength (dimensionless) \\
$\Lambda_{\max}$ & 1 & spectral radius of $\Lap$ (or sing.\ value) \\
$\Delta$ & 1 & max graph degree \\
$M\in\mathrm{GL}(2,\mathbb{C})$ & --- & transfer step; $V=\mathrm{tr}M$, $Q=\det M$ \\
$V_\ell, Q_\ell$ & --- & stride-$\ell$ invariants, $\mathrm{tr}(M^\ell)$, $(\det M)^\ell$ \\
$r_\pm=\rho e^{\pm i\theta}$ & --- & engine/AR roots \\
$\varepsilon(\mathbf q)$ & rad/step & band phase (Fourier) \\
$\Thet, F, \omega, \Omega$ & mixed & coframe, antisym.\ defect, connection, curvature \\
$h,\ell$ & length & mesh size, loop diameter \\
$\mu$ & $L^{-2}$ & spectral gap of $\Lap^\dagger\Lap$ \\
$D$ & 1/step & dephasing rate ($\tfrac12 S_\delta(0)$) \\
\bottomrule
\end{tabular}
\end{center}

% =======================================================================
\appendix
\section{Proof snippets and derivations}
\subsection*{A. Cayley--Hamilton $\Rightarrow$ decimated law}
Apply $\chi_\ell(M^\ell)=0$ to $M^\ell$; project as in \eqref{eq:engine-law}.
\subsection*{B. Finite-ring $G^{(L)}$ under net-zero forcing}
Fourier-diagonalize the ring operator to obtain \eqref{eq:ring-G}. Under a net-zero source the $k=0$ (DC) bin is absent (no DC response); handle the $V=2\sqrt{Q}$ resonance by the double-root limit/pseudoinverse.
\subsection*{C. Strang+BCH and leakage term}
Write the edge step $U_e=\E^{1/2} T \E^{1/2}$ with $T=e^{\tau\widehat K}$; BCH yields nested-commutator norms $\propto s$; concatenation over a loop plus the gap $\mu$ controls leakage as $\mu^{-1} s\,\ell$.

\section{Reproducibility: compact code cells (language-agnostic)}
\subsection*{AR(2) fit and dephasing estimator}
\begin{verbatim}
# Given complex residue y[0..N-1]
# Fit AR(2): y_{t+2} = V y_{t+1} - Q y_t + u_t
Vhat,Qhat = argmin_VQ sum_t |y[t+2]-V*y[t+1]+Q*y[t]|^2
rplus = 0.5*(Vhat + sqrt(Vhat**2 - 4*Qhat))
# Demodulate and phase increments
z[t] = y[t] * conj(rplus)**t
delta[t] = wrap(angle(z[t+1]*conj(z[t])) - angle(rplus))
# Bartlett-at-zero for S_delta(0); D = 0.5*S_delta(0)
\end{verbatim}

\section{Notes on operational vs intrinsic claims}
\emph{Intrinsic} quantities (e.g.\ $V_L=\mathrm{tr}(M^L)$) belong to the system and are invariant under clock twists; \emph{operational} estimators (e.g.\ $\widehat V_L$ from conditioned residues) \emph{can} change under PBC/APBC as the probe modifies the clock sector. Figures and statements are labeled accordingly.

\end{document}
